\section{Common Misconceptions and Mistakes}
    \subsection{Chem Core}
        \begin{itemize}
            \item Mass $\neq$ Weight — Mass is a measure of the amount of stuff, weight is a measure of gravitational expansion.
            \item Mass $\neq$ Density — Mass is a measure of the amount of stuff, Density is the amount of stuff in a given volume.
            \item Metric Prefixes and Conversions.
            \begin{center}
                \begin{tabular}{| c | c | c |}
                    \hline
                    Factor & Prefix & Symbol \\
                    \hline
                    $10^9$ & giga & G \\
                    $10^6$ & mega & M \\
                    $10^3$ & kilo & k \\
                    $10^{-2}$ & centi & c \\
                    $10^{-3}$ & milli & m\\
                    $10^{-6}$ & micro & $\mu$ \\
                    $10^{-9}$ & nano & n \\
                    $10^{-12}$ & pico & p\\
                    \hline
                \end{tabular}
            \end{center}
            Prefix $=$ Base $\times$ Factor (1 kg $=$ 1 g $\times 10^3$)
            \item Accuracy vs Precision. An measurement is accurate if it is close to the actual or correct value. A measurement is precise if repeated measurements are close together.
            \item Tools $\rightarrow$ Decimal $\rightarrow$ SigDigs $\rightarrow$ Answer
            \item Finding Exact measurements
            \item Rounding
            \item Dimensional Analysis
            \item amu (atomic mass unit) are very small and used to measure the mass of a singular atom. 
            \item Empirical Formula vs Molecular Formula. The empirical formula for a molecule represents the ratio of atoms in a molecule where a molecular formula represents the number of atoms. 
            \item Polyatomic ions do not disassociate in water
            \item Ionic compounds have nonmetal parts (\ce{NH4+},\ce{OH-}, ect)
            \item Acid names replace Ionic Names for Acids (Hydrochloric Acid not Hydrogen Chloride)
            \item Not all compounds beginning with hydrogen are acids (\ce{H2O},\ce{H2O2}, many organic molecules, ect)
            \item Molecular Prefixes (Mono, Di, Tri, ect) for Molecular compounds only
            \item Knowing how to name all compound types. 
            \item Coefficents $\neq$ Subscripts. Coefficients represent the number of molecules, while subscripts represent the number of atoms in a molecule.
            \item 1 mole = 6.022 $\times$ 10$^{23}$ particles. A mole is a number like a dozen, but instead of 12 it is 6.022 $\times$ 10$^{23}$, which is called Avogadro's number. It is used to count particles in chemistry.
            \item Mass $\neq$ Moles. Mass is a measure of the amount of stuff, moles is a measure of the number of particles. To convert between them, you need to use the molar mass.
            \item Moles in Lab $\neq$ moles in calculations. One is a theoretical value and one is a real value, they will not always be the same due to errors and inaccuracies in the lab.
            \item Coefficients in an empirical formula may have decimals. 
            \item Experimental Error must be identified.
            \item Reactant yielding the smallest product is the limiting reactant determining the theoretical yield.
            \item Percent yield can be greater than 100\% due to errors and impurities in the product.
            \item Show Work. You need to for the AP exam so start practicing now.
        \end{itemize}
    \subsection{Reactions}
    \begin{itemize}
        \item M = mol/L, not mols. Molarity is a measure of concentration, not amount.
        \item Atoms of the same element may have different oxidation numbers. 
        \item Oxidation does not need oxygen
        \item Strong acids and bases net ionic equations do have $\ce{H+_{(aq)} + OH-_{(aq)} -> H2O_{(l)}}$
        \item Equivalence Point and End point are not the same. Equivalence point is an ideal point based on math, while end point is the actual point in the lab that may be close to the equivalence point but is not exact.
        \item Non-electrolytes can still make ions.
        \item Pure water is a poor conductor
        \item Disolution and Dissociation are not the same. Dissolution is the process of a solute dissolving in a solvent, while dissociation is the process of an ionic compound breaking apart into its ions in solution.
        \item Ammonium (\ce{NH4+}) is a cation even though it is not a metal.
        \item Polyatomic ions do not disassociate in water, they stay together as a unit.
        \item Insolubility really means poorly soluble, not completely insoluble.
    \end{itemize}
    \subsection{Kinetics}
        \begin{itemize}
            \item Assuming reaction order from coefficients in the balanced equation is not always correct. 
            \item Zero Order processes
            \item Reaction rate from experimental results.
            \item Confuse fast reaction rate with large yields.
            \item Distinguish between kinetics vs. thermodynamics control of reaction rate. 
            \item Confuse intermediates with transition states. Intermediates are species that are formed during the reaction and can be isolated, while transition states are high-energy states that occur during the reaction and cannot be isolated.
            \item Confuse adsorption with absorption. Adsorption is the process of a substance sticking to the surface of another substance, while absorption is the process of a substance being taken up into the bulk of another substance.
        \end{itemize}
    \subsection{Equilibrium}
        \begin{itemize}
            \item Solids and liquids do not appear in the equilibrium constant expression because their concentrations do not change during the reaction.
            \item The values $k$ and $Q$ are different. $k$ is the equilibrium constant, which is a constant value for a given reaction at a given temperature, while $Q$ is the reaction quotient, which can change as the reaction progresses and is used to determine the direction in which the reaction will proceed to reach equilibrium.
            \item $k$ and $Q$ are unitless. They are relative ratios of concentrations and do not have units.
            \item The difference between $k_{eq}$, $k_a$, $k_b$, and $k_{sp}$. $k_{eq}$ is the equilibrium constant for a general reaction, $k_a$ is the acid dissociation constant, $k_b$ is the base dissociation constant, and $k_{sp}$ is the solubility product constant. They all represent different types of equilibrium constants for different types of reactions.
            \item Use the wrong arrow. Resonance arrow $\leftrightarrow$ is used for resonance structures, while equilibrium arrow $\rightleftharpoons$ is used for chemical equilibrium.
            \item Forgeting to include ``$K_{eq}=$'' in the final answer.
        \end{itemize}
    \subsection{Aqueous}
        \begin{itemize}
            \item Weak acids are different from low concentrations of strong acids. Weak means that it has an incomplete dissociation in water, while low concentration means that there is a small amount of a strong acid that is fully dissociated.
            \item Calculator use with scientific notation instead of standard notation
            \item Math with logarithms and exponential
            \item Confuse base 10 and natural logarithms
            \item Acids can be neutral or charged.
            \item $[H+]\times[OH-]=10^{-14}$ and $\text{pH} + \text{pOH} =14$ only at 25\textdegree C
            \item Neutralization reactions do not always end at a pH of 7
            \item Hydrogen ion, Hydronium Ion, Proton, and Aqueous Proton are all essentially equivalent and equal to $[H+]$. 
            \item Titration is not a reaction type but a laboratory procedure. 
            \item Everything about Buffers. Just go look at the buffers section.
            \item Volume changes when adding solutions together.
            \item Equilibrium affected by other reactions and concentrations in a solution. 
            \item Henderson Hasselbalch equation is only an approximation.
            \item There is a difference between $K_{sp}$ and Molar solubility.
            \item There is a difference between amphoteric and amphiprotic.
        \end{itemize}
    \subsection{Thermodynamics}
        \begin{itemize}
            \item Confuse Energy, Heat, Power, and Tempreture. Energy is the capacity to do work, heat is the transfer of energy due to a temperature difference, power is the rate at which work is done or energy is transferred, and temperature is a measure of the average kinetic energy of particles in a substance.
            \item An open container will be at a constant pressure.
            \item Enthalpy of reaction is different from molar enthalpy. Enthalpy of reaction is the total enthalpy change for a given reaction, while molar enthalpy is the enthalpy change per mole of a specific substance in the reaction.
            \item Formation reactions make exactly 1 mole of a compound from its elements in their standard states.
            \item Spontaneous reactions do not always occur quickly, and non-spontaneous reactions can occur if energy is added to the system.
            \item Change of thermodynamic values and their absolute values. 
            \item Signs for losses or gains of energy or temperature.
            \item Forget surroundings affect systems.
            \item Forget phases when calculating thermodynamic values.
            \item Standard conditions use 298.15 K. 
        \end{itemize}
    \subsection{Electrochemistry}
        \begin{itemize}
            \item Redox reactions.
            \item Redox reactions in acidic or basic solutions allow for the addition of $\ce{H+}$, $\ce{OH-}$, and $\ce{H2O}$ to balance the reaction.
            \item Oxidation does not need oxygen.
            \item In an electrochemical cell, electrons flow from the anode to the cathode.
            \item Individual electrode potentials cannot be measured, only the potential difference between two electrodes can be measured.
            \item Voltaic cells can be negative
            \item Cell potentials are affected by concentration of ions and gases.
            \item Electrochemical potentials are relative
            \item Electrochemical reactions can be reversible.
        \end{itemize}
    \subsection{AP Atom}
        \begin{itemize}
            \item Problems with conversions between nm, m, and Å.
            \item Difference between Certainty and Possibility. Just because something is possible does not mean it is certain or even likely.
            \item Don't know the basics of procedures such as spectroscopy.
            \item Confuse Böhr Orbits with Quantum Mechanical Orbitals. Böhr orbits are circular paths that electrons were once thought to follow around the nucleus, while quantum mechanical orbitals are regions of space where there is a high probability of finding an electron.
            \item Spectral lines are not the same as energy levels. Spectral lines are the specific wavelengths of light emitted or absorbed by an atom when electrons transition between energy levels, while energy levels are the specific energies that electrons can have in an atom.
            \item Spectral lines are difference between quantum energy levels, not the energy levels themselves. The energy of a spectral line corresponds to the difference in energy between two quantum energy levels, not the energy of the levels themselves.
            \item Orbital notation contain 2 electrons per orbital, not 1. Each orbital can hold a maximum of 2 electrons with opposite spins, so orbital notation should reflect this by showing pairs of electrons in each orbital.
            \item Difference between $s$, $p$, $d$, $f$ and $l$, $m_l$, $m_s$. The $s$, $p$, $d$, $f$ are the types of orbitals that describe the shape of the region where electrons are likely to be found, while $l$, $m_l$, and $m_s$ are quantum numbers that describe the energy level, orientation, and spin of electrons in those orbitals.
            Problems with electron configurations.
            \item Shielding between electrons does not increase going left to right across a period ($\rightarrow$)
            \item Effective nuclear charge increase going left to right across a period ($\rightarrow$)
            \item Ionic charge is different than nuclear charge
            \item Atomic radius decrease going left to right across a period ($\rightarrow$)
            Signs of electron affinities.
            \item Behavior of aqueous elements with set phases.
            \item Ionic Radius refers to solids, Ionization Energy and Electron Affinity refer to gaseous atoms.
            \item Isoelectronic does not mean that the atoms have the same number of valence electrons.
            \item Triple bonds are not three times strong than single bonds. Double and triple bonds are stronger than single bonds, but the strength does not increase linearly with the number of bonds. The strength of a bond depends on various factors such as the types of atoms involved, the bond length, and the overall molecular structure.
            \item The first bond formed in a molecule, the $\sigma$ bond is the strongest bond, while the second and third bonds, the $\pi$ bonds, are weaker. In a double bond, there is one $\sigma$ bond and one $\pi$ bond, while in a triple bond, there is one $\sigma$ bond and two $\pi$ bonds. The $\sigma$ bond is stronger than the $\pi$ bonds because it is formed by the head-on overlap of orbitals, while the $\pi$ bonds are formed by the side-to-side overlap of orbitals, which is less effective in creating a strong bond.
            \item Resonance arrow $\leftrightarrow$ is not an equilibrium arrow $\rightleftharpoons$.
            \item Resonance shows delocalization of electrons, which means that there are many possible Lewis structures.
            \item There are many exceptions to the octet rule.
            \item Eight valence electrons does not always mean that the atom is stable. 
            \item Formal charge is different from ionic/real charge, which is different from effective nuclear charge. Formal charge is a theoretical charge assigned to an atom in a molecule based on the assumption that electrons are shared equally in bonds, while ionic/real charge is the actual charge on an ion or atom based on the number of electrons it has gained or lost, and effective nuclear charge is the net positive charge experienced by an electron in an atom, taking into account both the actual nuclear charge and the shielding effect of other electrons.
            \item Not all polar substances conduct electricity. Take water for example.
            \item Lewis Dots are the best attempt at a 2D representation of a 3D molecule. 
            \item Three dimensional space is hard to visualize.
            \item Electron domain geometry is not the same as molecular geometry. Electron domain geometry considers the arrangement of all electron domains (bonding and non-bonding) around the central atom, while molecular geometry only considers the arrangement of atoms (bonding domains) and ignores lone pairs.
            \item Large molecules are very hard to work with so work on smaller parts and put them together. 
            \item Finding the dipole moment of a molecule requires finding the molecular geometry of the molecule.
            \item The presence of polar bonds does not necessarily mean that the molecule is polar since the direction of the polar bonds and the molecular geometry can result in the individual dipole moments canceling each other out, resulting in a nonpolar molecule. 
            \item Ions have no polarity since they are separate charged particles and electrons are not shared between them.
            \item Hybridization is connected to the electron domain not the molecular geometry. 
            \item Bonds can interfere with each other in constructive ways.
            \end{itemize}
    \subsection{States of Matter}
            \begin{itemize}
                \item Always use Kelvin for gas temperatures.
                \item Tracking units is especially important in gas law problems, iuncluding the ideal gas constant $R$.
                \item STP and standard conditions are different. STP is 0\textdegree C and 1 atm, while standard conditions is 25\textdegree C and 1 atm. We use STP for gas law problems and standard conditions for thermodynamics problems.
                \item Ideal gases are just an approximation and real gases deviate from reality. This is especially true at high pressures and low temperatures.
                \item There is a difference between effusion and diffusion. Effusion is the process of gas particles passing through a tiny opening, while diffusion is the process of gas particles spreading out and mixing with other gases.
                \item Sum of mole fraction must equal 1. 
                \item Gas densities are commonly expressed in g/L, not g/mL. 
                \item Differences between intermolecular forces and intramolecular forces. Intermolecular forces are the forces of attraction or repulsion between molecules, while intramolecular forces are the forces that hold atoms together within a molecule.
                \item The relative strength of intermolecular forces.
                \item Unaware of intramolecular hydrogen bonding. 
                \item Adhesion and Cohesion are different. Adhesion is the attraction between different substances, while cohesion is the attraction between particles of the same substance.
                \item Water can sublime under the right conditions.
                \item Viscous does not mean more dense. 
                \item Liquid crystals are liquids, but also crystalline.
                \item Ion-dipole forces are technically Interparticle forces, but we pretend they are intermolecular forces.
                \item Boiling and melting points depend on atmospheric pressure, not just the substance itself.
                \item Distinguishing between unit cells.
                \item Solutions are not always aqueous or water-based.
                \item Alloys are solid solutions of metals.
                \item Foams are liquids or solids that are dissolved in a gas.
                \item Ionic Solids attractions based on the charges or strength of the ions.
                \item Nanomaterials are not all man-made.
                \item Polymers are repeating covalent segments of atoms.
                \item Determining the properties of non-crystalline solids.
                \item Distinguishing the difference bteween the types of solids
                \item Shapes of crystal lattices is different from the shapes of VSEPR molecules. 
                \item Confusion between weak and strong compared to dilute and concentrated. Weak and strong refer to the degree of dissociation of an acid or base, while dilute and concentrated refer to the amount of solute in a given volume of solution.
                \item Fail to identify the driving forces for solution formation
                \item Confusion between dissolution and melting
                \item Crystallization is the reverse of dissolution. 
                \item There are nine phase combinations for solutions, not just solids in liquids. 
                \item A solution is a uniform mixture where solvation is the process of solute solvent interactions.
                \item Not every mixture is a solution.
                \item Molality uses the mass of the solvent not the mass of the solution.
                \item Colloids can be in all three phases, not just liquids.
                \item Even insoluble substances can dissolve a little bit in water, they are just poorly soluble.
            \end{itemize}